% P10-Core-Inhaltsprüfung, 15.09.2026: datierter Ergebnisanschluss; B-60 korrigiert.
% Belege: thesis-evidence/w2/core-content-audit-20260915/evaluation/
% B-58, 15.09.2026: 201 -> 207 -> 208 -> 212 Items, keine Löschung im geprüften Abschnitt.
% Nachtrag: thesis-evidence/w2/f1-zaehlkorrektur-20260915.md; historische Runs unverändert.
% ============================================================
% Tabelle 7.5-1: Ergebnisübersicht der Fallserie — v01 (07.09.)
% In Overleaf ablegen als: tables/Chap7-v02/chap7.5table1.tex
% Aufruf im Fließtext: % P10-Core-Inhaltsprüfung, 15.09.2026: datierter Ergebnisanschluss; B-60 korrigiert.
% Belege: thesis-evidence/w2/core-content-audit-20260915/evaluation/
% B-58, 15.09.2026: 201 -> 207 -> 208 -> 212 Items, keine Löschung im geprüften Abschnitt.
% Nachtrag: thesis-evidence/w2/f1-zaehlkorrektur-20260915.md; historische Runs unverändert.
% ============================================================
% Tabelle 7.5-1: Ergebnisübersicht der Fallserie — v01 (07.09.)
% In Overleaf ablegen als: tables/Chap7-v02/chap7.5table1.tex
% Aufruf im Fließtext: % P10-Core-Inhaltsprüfung, 15.09.2026: datierter Ergebnisanschluss; B-60 korrigiert.
% Belege: thesis-evidence/w2/core-content-audit-20260915/evaluation/
% B-58, 15.09.2026: 201 -> 207 -> 208 -> 212 Items, keine Löschung im geprüften Abschnitt.
% Nachtrag: thesis-evidence/w2/f1-zaehlkorrektur-20260915.md; historische Runs unverändert.
% ============================================================
% Tabelle 7.5-1: Ergebnisübersicht der Fallserie — v01 (07.09.)
% In Overleaf ablegen als: tables/Chap7-v02/chap7.5table1.tex
% Aufruf im Fließtext: % P10-Core-Inhaltsprüfung, 15.09.2026: datierter Ergebnisanschluss; B-60 korrigiert.
% Belege: thesis-evidence/w2/core-content-audit-20260915/evaluation/
% B-58, 15.09.2026: 201 -> 207 -> 208 -> 212 Items, keine Löschung im geprüften Abschnitt.
% Nachtrag: thesis-evidence/w2/f1-zaehlkorrektur-20260915.md; historische Runs unverändert.
% ============================================================
% Tabelle 7.5-1: Ergebnisübersicht der Fallserie — v01 (07.09.)
% In Overleaf ablegen als: tables/Chap7-v02/chap7.5table1.tex
% Aufruf im Fließtext: \input{tables/Chap7-v02/chap7.5table1}
% Label: t:eval-fallserie
% Stil: chap6.1-Konvention.
% HERKUNFT: thesis-evidence/w2/z2-fallblaetter.md rev2
% (Ergebnistabelle; Urteils-Legende dort im Kopf). Läufe je Zeile
% siehe chap7.5.tex-HERKUNFT.
% ============================================================

\begin{table}[htbp]
  \centering
  \small
  \begin{tabular}{>{\raggedright\arraybackslash}p{0.16\textwidth}>{\raggedright\arraybackslash}p{0.12\textwidth}llll}
    \hline
    \textbf{Falltyp} & \textbf{Lauf (Stand)} & \textbf{Inhalt} & \textbf{Freigabe} & \textbf{Herkunft} & \textbf{Projektion} \\
    \hline
    Neue Anforderung & 18.08. & erfüllt\textsuperscript{a} & erfüllt\textsuperscript{a} & erfüllt & erfüllt (real) \\
    Präzisierung & 17.08. & teilweise\textsuperscript{f} & erfüllt & erfüllt & Dry-Run \\
    Widerspruch & 04.08. & erfüllt\textsuperscript{b} & erfüllt\textsuperscript{c} & \textbf{abweichend} & nicht untersucht \\
    Wiederholung & 18.08. & n.\,b. & erfüllt & erfüllt & n.\,a. \\
    Offene Frage & 18.08. & erfüllt & erfüllt & erfüllt & n.\,a.\textsuperscript{d} \\
    Ext.\ Rückweg & 20.08. & erfüllt & erfüllt & erfüllt & Hauptkette erfüllt\textsuperscript{e} \\
    \hline
    \multicolumn{6}{p{0.95\textwidth}}{\footnotesize
      Die Einordnung wurde fallbezogen beurteilt; bei F2 ist die
      Rücknahme des Kalender-Prüfauftrags nicht gesondert begründet
      (Fußnote~f). \textsuperscript{a}~für die verfolgten
      Ketten/belegten Operationen, kein Pauschalurteil über alle
      11 vor der GitHub-Projektion neu hinzugekommenen Core-Items.
\textsuperscript{b}~Ablösung und Relationen;
      fachliche Angleichung der abhängigen Backlog-Items blieb
      offener Klärungsbedarf. \textsuperscript{c}~nur das
      Entscheidungs-Gate lief interaktiv, die übrigen vier Gates im
      deklarierten Experiment-Modus. \textsuperscript{d}~offene
      Entscheidungen haben vertragsgemäß keinen Projektionskanal.
      \textsuperscript{e}~ein Vermerk-Schreibvorgang zunächst
      fehlgeschlagen; Run-Zuordnung des Nachtrags n.\,b.
      \textsuperscript{f}~Nachprüfung vom 15.09.: neue Vorgaben
      erhalten, alter Kalender-Prüfauftrag nur historisiert;
      Rücknahme nicht gesondert begründet. Dry-Run = Projektion
      vorbereitet, nicht ausgeführt.
      Legende: n.\,b. = Beleg fehlt · n.\,a. = vertragsgemäß nicht
      vorgesehen · „nicht untersucht`` = Prüfumfang bewusst
      begrenzt --- keine dieser Angaben ist eine Erfüllung.} \\
    \hline
  \end{tabular}
  \caption[Ergebnisse der fachlichen Fallserie]{Ergebnisübersicht der fachlichen Fallserie (fünf
    Falltypen, historisch belegte Läufe mit je eigenem
    Systemstand; keine Erfolgsquote). Beobachtete Abweichungen und
    Nachweislücken bleiben getrennt ausgewiesen: die
    Herkunfts-\emph{Anzeige} im Widerspruchsfall (abweichend), der
    zunächst fehlgeschlagene Vermerk-Schreibvorgang im Rückweg-Fall
    (Fußnote~e) sowie die n.\,b.-Zellen.}
  \label{t:eval-fallserie}
\end{table}

% Label: t:eval-fallserie
% Stil: chap6.1-Konvention.
% HERKUNFT: thesis-evidence/w2/z2-fallblaetter.md rev2
% (Ergebnistabelle; Urteils-Legende dort im Kopf). Läufe je Zeile
% siehe chap7.5.tex-HERKUNFT.
% ============================================================

\begin{table}[htbp]
  \centering
  \small
  \begin{tabular}{>{\raggedright\arraybackslash}p{0.16\textwidth}>{\raggedright\arraybackslash}p{0.12\textwidth}llll}
    \hline
    \textbf{Falltyp} & \textbf{Lauf (Stand)} & \textbf{Inhalt} & \textbf{Freigabe} & \textbf{Herkunft} & \textbf{Projektion} \\
    \hline
    Neue Anforderung & 18.08. & erfüllt\textsuperscript{a} & erfüllt\textsuperscript{a} & erfüllt & erfüllt (real) \\
    Präzisierung & 17.08. & teilweise\textsuperscript{f} & erfüllt & erfüllt & Dry-Run \\
    Widerspruch & 04.08. & erfüllt\textsuperscript{b} & erfüllt\textsuperscript{c} & \textbf{abweichend} & nicht untersucht \\
    Wiederholung & 18.08. & n.\,b. & erfüllt & erfüllt & n.\,a. \\
    Offene Frage & 18.08. & erfüllt & erfüllt & erfüllt & n.\,a.\textsuperscript{d} \\
    Ext.\ Rückweg & 20.08. & erfüllt & erfüllt & erfüllt & Hauptkette erfüllt\textsuperscript{e} \\
    \hline
    \multicolumn{6}{p{0.95\textwidth}}{\footnotesize
      Die Einordnung wurde fallbezogen beurteilt; bei F2 ist die
      Rücknahme des Kalender-Prüfauftrags nicht gesondert begründet
      (Fußnote~f). \textsuperscript{a}~für die verfolgten
      Ketten/belegten Operationen, kein Pauschalurteil über alle
      11 vor der GitHub-Projektion neu hinzugekommenen Core-Items.
\textsuperscript{b}~Ablösung und Relationen;
      fachliche Angleichung der abhängigen Backlog-Items blieb
      offener Klärungsbedarf. \textsuperscript{c}~nur das
      Entscheidungs-Gate lief interaktiv, die übrigen vier Gates im
      deklarierten Experiment-Modus. \textsuperscript{d}~offene
      Entscheidungen haben vertragsgemäß keinen Projektionskanal.
      \textsuperscript{e}~ein Vermerk-Schreibvorgang zunächst
      fehlgeschlagen; Run-Zuordnung des Nachtrags n.\,b.
      \textsuperscript{f}~Nachprüfung vom 15.09.: neue Vorgaben
      erhalten, alter Kalender-Prüfauftrag nur historisiert;
      Rücknahme nicht gesondert begründet. Dry-Run = Projektion
      vorbereitet, nicht ausgeführt.
      Legende: n.\,b. = Beleg fehlt · n.\,a. = vertragsgemäß nicht
      vorgesehen · „nicht untersucht`` = Prüfumfang bewusst
      begrenzt --- keine dieser Angaben ist eine Erfüllung.} \\
    \hline
  \end{tabular}
  \caption[Ergebnisse der fachlichen Fallserie]{Ergebnisübersicht der fachlichen Fallserie (fünf
    Falltypen, historisch belegte Läufe mit je eigenem
    Systemstand; keine Erfolgsquote). Beobachtete Abweichungen und
    Nachweislücken bleiben getrennt ausgewiesen: die
    Herkunfts-\emph{Anzeige} im Widerspruchsfall (abweichend), der
    zunächst fehlgeschlagene Vermerk-Schreibvorgang im Rückweg-Fall
    (Fußnote~e) sowie die n.\,b.-Zellen.}
  \label{t:eval-fallserie}
\end{table}

% Label: t:eval-fallserie
% Stil: chap6.1-Konvention.
% HERKUNFT: thesis-evidence/w2/z2-fallblaetter.md rev2
% (Ergebnistabelle; Urteils-Legende dort im Kopf). Läufe je Zeile
% siehe chap7.5.tex-HERKUNFT.
% ============================================================

\begin{table}[htbp]
  \centering
  \small
  \begin{tabular}{>{\raggedright\arraybackslash}p{0.16\textwidth}>{\raggedright\arraybackslash}p{0.12\textwidth}llll}
    \hline
    \textbf{Falltyp} & \textbf{Lauf (Stand)} & \textbf{Inhalt} & \textbf{Freigabe} & \textbf{Herkunft} & \textbf{Projektion} \\
    \hline
    Neue Anforderung & 18.08. & erfüllt\textsuperscript{a} & erfüllt\textsuperscript{a} & erfüllt & erfüllt (real) \\
    Präzisierung & 17.08. & teilweise\textsuperscript{f} & erfüllt & erfüllt & Dry-Run \\
    Widerspruch & 04.08. & erfüllt\textsuperscript{b} & erfüllt\textsuperscript{c} & \textbf{abweichend} & nicht untersucht \\
    Wiederholung & 18.08. & n.\,b. & erfüllt & erfüllt & n.\,a. \\
    Offene Frage & 18.08. & erfüllt & erfüllt & erfüllt & n.\,a.\textsuperscript{d} \\
    Ext.\ Rückweg & 20.08. & erfüllt & erfüllt & erfüllt & Hauptkette erfüllt\textsuperscript{e} \\
    \hline
    \multicolumn{6}{p{0.95\textwidth}}{\footnotesize
      Die Einordnung wurde fallbezogen beurteilt; bei F2 ist die
      Rücknahme des Kalender-Prüfauftrags nicht gesondert begründet
      (Fußnote~f). \textsuperscript{a}~für die verfolgten
      Ketten/belegten Operationen, kein Pauschalurteil über alle
      11 vor der GitHub-Projektion neu hinzugekommenen Core-Items.
\textsuperscript{b}~Ablösung und Relationen;
      fachliche Angleichung der abhängigen Backlog-Items blieb
      offener Klärungsbedarf. \textsuperscript{c}~nur das
      Entscheidungs-Gate lief interaktiv, die übrigen vier Gates im
      deklarierten Experiment-Modus. \textsuperscript{d}~offene
      Entscheidungen haben vertragsgemäß keinen Projektionskanal.
      \textsuperscript{e}~ein Vermerk-Schreibvorgang zunächst
      fehlgeschlagen; Run-Zuordnung des Nachtrags n.\,b.
      \textsuperscript{f}~Nachprüfung vom 15.09.: neue Vorgaben
      erhalten, alter Kalender-Prüfauftrag nur historisiert;
      Rücknahme nicht gesondert begründet. Dry-Run = Projektion
      vorbereitet, nicht ausgeführt.
      Legende: n.\,b. = Beleg fehlt · n.\,a. = vertragsgemäß nicht
      vorgesehen · „nicht untersucht`` = Prüfumfang bewusst
      begrenzt --- keine dieser Angaben ist eine Erfüllung.} \\
    \hline
  \end{tabular}
  \caption[Ergebnisse der fachlichen Fallserie]{Ergebnisübersicht der fachlichen Fallserie (fünf
    Falltypen, historisch belegte Läufe mit je eigenem
    Systemstand; keine Erfolgsquote). Beobachtete Abweichungen und
    Nachweislücken bleiben getrennt ausgewiesen: die
    Herkunfts-\emph{Anzeige} im Widerspruchsfall (abweichend), der
    zunächst fehlgeschlagene Vermerk-Schreibvorgang im Rückweg-Fall
    (Fußnote~e) sowie die n.\,b.-Zellen.}
  \label{t:eval-fallserie}
\end{table}

% Label: t:eval-fallserie
% Stil: chap6.1-Konvention.
% HERKUNFT: thesis-evidence/w2/z2-fallblaetter.md rev2
% (Ergebnistabelle; Urteils-Legende dort im Kopf). Läufe je Zeile
% siehe chap7.5.tex-HERKUNFT.
% ============================================================

\begin{table}[htbp]
  \centering
  \small
  \begin{tabular}{>{\raggedright\arraybackslash}p{0.16\textwidth}>{\raggedright\arraybackslash}p{0.12\textwidth}llll}
    \hline
    \textbf{Falltyp} & \textbf{Lauf (Stand)} & \textbf{Inhalt} & \textbf{Freigabe} & \textbf{Herkunft} & \textbf{Projektion} \\
    \hline
    Neue Anforderung & 18.08. & erfüllt\textsuperscript{a} & erfüllt\textsuperscript{a} & erfüllt & erfüllt (real) \\
    Präzisierung & 17.08. & teilweise\textsuperscript{f} & erfüllt & erfüllt & Dry-Run \\
    Widerspruch & 04.08. & erfüllt\textsuperscript{b} & erfüllt\textsuperscript{c} & \textbf{abweichend} & nicht untersucht \\
    Wiederholung & 18.08. & n.\,b. & erfüllt & erfüllt & n.\,a. \\
    Offene Frage & 18.08. & erfüllt & erfüllt & erfüllt & n.\,a.\textsuperscript{d} \\
    Ext.\ Rückweg & 20.08. & erfüllt & erfüllt & erfüllt & Hauptkette erfüllt\textsuperscript{e} \\
    \hline
    \multicolumn{6}{p{0.95\textwidth}}{\footnotesize
      Die Einordnung wurde fallbezogen beurteilt; bei F2 ist die
      Rücknahme des Kalender-Prüfauftrags nicht gesondert begründet
      (Fußnote~f). \textsuperscript{a}~für die verfolgten
      Ketten/belegten Operationen, kein Pauschalurteil über alle
      11 vor der GitHub-Projektion neu hinzugekommenen Core-Items.
\textsuperscript{b}~Ablösung und Relationen;
      fachliche Angleichung der abhängigen Backlog-Items blieb
      offener Klärungsbedarf. \textsuperscript{c}~nur das
      Entscheidungs-Gate lief interaktiv, die übrigen vier Gates im
      deklarierten Experiment-Modus. \textsuperscript{d}~offene
      Entscheidungen haben vertragsgemäß keinen Projektionskanal.
      \textsuperscript{e}~ein Vermerk-Schreibvorgang zunächst
      fehlgeschlagen; Run-Zuordnung des Nachtrags n.\,b.
      \textsuperscript{f}~Nachprüfung vom 15.09.: neue Vorgaben
      erhalten, alter Kalender-Prüfauftrag nur historisiert;
      Rücknahme nicht gesondert begründet. Dry-Run = Projektion
      vorbereitet, nicht ausgeführt.
      Legende: n.\,b. = Beleg fehlt · n.\,a. = vertragsgemäß nicht
      vorgesehen · „nicht untersucht`` = Prüfumfang bewusst
      begrenzt --- keine dieser Angaben ist eine Erfüllung.} \\
    \hline
  \end{tabular}
  \caption[Ergebnisse der fachlichen Fallserie]{Ergebnisübersicht der fachlichen Fallserie (fünf
    Falltypen, historisch belegte Läufe mit je eigenem
    Systemstand; keine Erfolgsquote). Beobachtete Abweichungen und
    Nachweislücken bleiben getrennt ausgewiesen: die
    Herkunfts-\emph{Anzeige} im Widerspruchsfall (abweichend), der
    zunächst fehlgeschlagene Vermerk-Schreibvorgang im Rückweg-Fall
    (Fußnote~e) sowie die n.\,b.-Zellen.}
  \label{t:eval-fallserie}
\end{table}
